\section{材料与方法}
\label{sec:methods}

\subsection{数据、窗口与经验目标}
\label{sec:meth-data}

使用 2026 年 10 月 20 日检索的美国国家医疗保健安全网络（NHSN）每周呼吸道住院数据快照，流感分析限于 50 州和哥伦比亚特区\citep{cdc2026nhsn,frutos2025evaluation}，排除全国、区域汇总及领地。医院参与率在观测期内变化，其中 2024 年 5--10 月为自愿报告期，之后恢复报告要求\citep{cms2024reporting}。本文描述该监测系统，不估计制度变动的因果效应。

主窗口保留此前披露的四个流感季早期增长阶段：2022 年 10 月 1 日至 11 月 26 日、2023 年 10 月 7 日至 12 月 23 日、2024 年 11 月 9 日至 12 月 28 日、2025 年 10 月 4 日至 12 月 27 日。每个窗口是该季允许起点后、第一个全国转折点之前，全国周增长不低于 1.05 的连续片段；2024 年起点在恢复强制报告后。窗口规则与主统计量已在先前分析轮次中选择，本轮修订不是预注册研究。我们固定这些窗口，而未重新选择有利于正截距的片段。推断以已选窗口为条件。

从原始快照重建观测，缺失或抑制值保持缺失，不按零住院处理；真实零计数保留。一行观测要求分母周之前两周至下一周的四个计数均存在，且日期严格按七天连续。加权病例规模需满足 $D_t\ge5$，主单周统计量还要求 $C_t\ge5$。得到 $R_A$ 的 \Nra{} 个州周与 $G$ 的 \Ngs{} 个州周，覆盖 \Nstates{} 个辖区、\Nweeks{} 个日历周。样本差异由分母条件决定。稀释之前固定州周键及原始计数。

\subsection{表观统计量与相对平方偏差}
\label{sec:meth-stats}

令 $C_t$ 表示第 $t$ 周住院数。加权增长统计量为
\begin{equation}
G_t=\frac{C_{t+1}}{D_t},\qquad D_t=0.65C_t+0.25C_{t-1}+0.10C_{t-2},
\label{eq:dt}
\end{equation}
增长率到再生数的变换为
\begin{equation}
R_{A,t}=\left[1+\frac{\theta}{7}\log\frac{C_{t+1}+0.5}{C_t+0.5}\right]_{+}^{\alpha},\qquad
\alpha=(g_\mu/g_\sigma)^2,\quad\theta=g_\sigma^2/g_\mu.
\label{eq:ra}
\end{equation}
正部约定使底数非正时仍有数值定义，并不意味着指数增长近似在此成立。伽马间隔映射来自增长率与再生数的更新关系\citep{wallinga2007how}。两个统计量都不是真实再生数的观测值，也不是感染更新模型给出的后验估计\citep{cori2013new,eales2024differences}。

为保持历史可比性，主变换采用 $g_\mu=3.6$ 天、$g_\sigma=1.6$ 天，即 \citet{cowling2009estimation} 的家庭序列间隔估计。这是代际间隔代理值。敏感性分别采用均值 3.2 天并固定原变异系数，以及均值 3.2 天、标准差 2.1 天；后者是 \citet{chan2025estimating} 表 2 的内禀代际间隔估计。固定变异系数是本文假设。家庭内实现代际间隔和亚型序列间隔是不同测量对象\citep{chan2025estimating,vink2014serial}。

每个设定计算
\begin{equation}
L_{st}=\left(\frac{X_{st}-\bar X}{\bar X}\right)^2,
\qquad E(L_{st}\mid D_{st})=a+b/D_{st},
\label{eq:affine}
\end{equation}
其中 $X$ 为表观统计量，$\bar X$ 为合并均值。对于 $R_A$，回归中的 $D_t$ 替换为 $C_t$。响应量衡量相对于拟合参照的离散程度，不是围绕已知真实 $R_t$ 的估计误差平方。$b/a$ 有定义时，描述该统计量病例尺度上的仿射工作模型，并不直接等于潜在平衡规模 $\mx$。

\subsection{拟合、不确定性与模型比较}
\label{sec:meth-inference}

采用迭代可行广义最小二乘（FGLS）拟合仿射均值，精度权重与拟合均值平方的倒数成比例。截距与斜率可以有正负号，但观测行上的拟合均值必须为正。阻尼更新使伽马准偏差目标 $\sum_iw_i[L_i/\mu_i+\log\mu_i]$ 下降；必要时的数值备用求解仍使用同一目标及正拟合均值约束。记录收敛性、秩与失败情况，达到迭代上限不作为有效估计。该目标是估计准则，不表示经验损失服从伽马分布。

主合并区间使用 2000 次州 $\times$ 日历周鸽巢自助抽样\citep{owen2007pigeonhole}。行权重为独立抽取的州与周多重性乘积，每次重新估计中心和回归；次要分析采用 1000 次抽样。另结合州多重性，在各季内按两周、三周连续块抽样，以检查全国共同冲击持续性的影响。四个季节不足以支持对任意未来流行机制的无约束外推。未定义比值、分母符号和拟合失败均保留诊断；若区间仅由有限有效重复计算，明确标注。

通过逐州留一与逐季留一比较仿射均值和 $\lambda D^{-\gamma}$。响应中心仅由训练观测计算，测试损失围绕训练中心构造；两个模型也仅在训练折内拟合。测试评分为 $L/\mu+\log\mu$ 的均值，并保留平方预测误差作为次要指标。报告幂律评分减去仿射评分，正值有利于仿射模型。训练集重叠的折不是独立重复，也不以仅四个流感季构造过度精确的显著性检验。此前残差自助平台检验的 $p$ 值不作为已修复检验继续沿用。

\subsection{受控二项稀释}
\label{sec:meth-thinning}

在 $q\in\{0.8,0.6,0.5,0.4,0.3,0.2,0.1\}$ 下，每个唯一州周计数每次从 $\operatorname{Bin}(C,q)$ 抽取一次，所有相邻比值复用该值；不重新应用筛选条件或窗口。重新计算统计量、中心和仿射拟合，回归尺度仍是原始 $C_t$。点响应平均 \ThinPointReps{} 次分别拟合的实现；\ThinBootstrap{} 次外层自助样本中，每水平平均 \ThinInnerReps{} 次内层拟合。若估计的蒙特卡洛标准误超过预先规定的样本不确定性比例，统一增加内层次数；报告实际预算与剩余模拟误差。

近似关系为
\begin{equation}
b(q)-b(1)\simeq c_{\rm thin}(1/q-1).
\label{eq:thinning}
\end{equation}
估计过原点的响应斜率，保留各水平，并分别剔除 $q=0.1$ 和逐一剔除其他水平。预定子样本要求原始 $C_t\ge50$，或同时要求 $C_t,C_{t+1}\ge50$；这限制的是稀释计数期望，而非每次实现。参照系数为 $\kappa(g)^2[1+(C_t+0.5)/(C_{t+1}+0.5)]$ 的加权均值，其中 $\kappa(g)=(g_\mu/7)/[1+\theta\log(g)/7]$、$g=(C_{t+1}+0.5)/(C_t+0.5)$。该表达式近似变换的相对响应，并稳定倒数计数，不是重新中心化后非线性统计量的精确方差。每次重抽样重新计算参照值，以形成配对 $c_{\rm thin}/b_P$ 区间；用同时区间带描述截距差向量。已知泊松计数例经过同一流程，仅用于机制检查，不作为区间覆盖率验证。

\subsection{州异质性与参照敏感性}
\label{sec:meth-fe}

保持同一合并损失定义，拟合 $E(L_{st}\mid C_{st})=a_s+b_{\rm w}/C_{st}$。州摘要以抽中的州出现次数为目标：重复抽中的州重复计入，未抽中的州不以零系数参加平均；行和病例加权摘要使用重抽样总量。合并与州内模型在同一自助样本中配对比较。分母敏感性单独把 $C_t$ 换成实际的加权 $D_t$。仅周和州加周模型作为数值诊断，失败拟合明确记为不可用。

另一组分析改变损失参照，分别使用季均值、各季内所有州共享的线性时间趋势、全国周均值，以及排除自身州的全国周均值。统一保持全局均值平方作为归一化分母，以区分中心改变与尺度改变；中心和趋势在自助样本中重估。这些操作改变经验目标，也可能移除真实传播信号，因此截距差是定义敏感性，不是正交生物学方差分解。另报告有符号相对偏差 $d=X/\bar X-1$ 的 $\operatorname{Var}_w(\bar d_w)/\operatorname{Var}_{st}(d_{st})$，明确称为描述性方差比，而不称平方损失的 ICC。

\subsection{条件校准与设计示例}
\label{sec:meth-inversion}

日度更新模拟器保持四个固定窗口长度及各州初始计数。传播采用独立负二项后代，周度再生数冲击区分全国和州分量，并声明其边际方差及持续性；住院延迟和报告概率除有外部依据者外均列为假设。二项住院事件按延迟分配；另一观察模型直接从潜在延迟住院强度 $\lambda$ 生成计数，条件均值为 $\lambda$、条件方差为 $\varphi\lambda$，不对已经带噪声的住院数再次抽样。$\varphi=1$ 为泊松极限，单独的观察矩测试检查该定义。

模拟面板使用与实证相同的统计量、筛选和拟合函数。方差参数描述该模拟器中的周度再生数冲击，不换算成每代方差或操作性周病例门槛。联合经验截距抽样和联合模拟面板重抽样传播样本与校准蒙特卡洛不确定性。不同代际间隔及观察模型情景表示结构敏感性，不是机制的后验分布，也不是覆盖全部机制的联合置信区间。补充方法列出全部设定、边界和模拟诊断；另报告分季、医院参与率敏感性以及明确假设下的曲线辨别设计实验，不将其用于报告制度因果效应或普适监测样本量推断。
